%%%%%%%%%%%%%%%%%%%%%%%%%%%%%%%%%%%%%%%%%%%%%%%%%%%%%%%%%%%%
\section[% Kurzversion für das Inhaltsverz., Kolumnentitel und PDF-Lesezeichen:
         Wichtiges zu Umbrüchen bei Überschriften
         \mbox{(Kurzversion für das Inhaltsverzeichnis \etc)}%
				]{% Langversion, die im Text gedruckt wird:
         Wichtiges zu Umbrüchen bei Überschriften
         (und ein Beispiel \newline für eine lange Überschrift,
         welche \newline für das Inhaltsverzeichnis und 
         \newline die Kolumnentiteln zu lang ist)%
				}%
\index{Zeilenumbruch}\index{Titel}\index{Ueberschrift@Überschrift}\index{PDF!Lesezeichen}%
\label{sec:Titles}
%%%%%%%%%%%%%%%%%%%%%%%%%%%%%%%%%%%%%%%%%%%%%%%%%%%%%%%%%%%%
%
%
Bei den Kapitelüberschriften kann man zwei Versionen definieren:
eine lange Überschrift in geschweiften Klammern, welche in der Arbeit selbst angezeigt wird, 
und optional eine Kurzversion in eckigen Klammern, welche im Inhaltsverzeichnis und in den 
\href{https://de.wikipedia.org/wiki/Kolumnentitel}{Kolumnentiteln}%
\footnote{Kolumnentitel sind Überschriften der einzelnen Seiten. Meist stehen sie in der Kopfzeile.}
angezeigt wird:

\begin{inlinelatex}
\section[Kurzversion]{Langer Titel}
\end{inlinelatex}

Dasselbe gilt für Bild- und Tabellenunterschriften. Hier kann man dem
\lc{caption}-Befehl ebenfalls einen optionalen Parameter in eckigen Klammern übergeben.

Um bei längeren Überschriften ein besseres Zeilenumbruchsverhalten im Inhaltsverzeichnis zu erzwingen
(\zB um eine ähnliche Länge aller Zeilen bei mehrzeiligen Überschriften zu erreichen),
ist manchmal eine Einfügung manueller Zeilenumbrüche in der Kurzversion der Überschriften erforderlich.
Allerdings funktioniert ein Einfügen manueller Zeilenumbrüche in der Kurzversion des Titels
etwa mit \verb+\newline+  leider nicht, da diese dann nicht nur im Inhaltsverzeichnis,
sondern auch in den Kolumnentiteln und PDF"=Lesezeichen zur Geltung kommen, 
was normalerweise nicht erwünscht ist.

Abhilfe schafft der folgende Trick:
man schließt den letzten, umzubrechenden Teil der Kurzversion des Titels in eine \verb+\mbox{}+.
Der Text, der in eine \verb+\mbox{}+ eingeschlossen wird, darf nicht umbrochen werden.
Im Kolumnentiteln und in den PDF"=Lesezeichen hat dies keine besondere Wirkung;
im Inhaltsverzeichnis führt dies jedoch dazu, dass \LaTeX{} den Zeilenumbruch vor der \verb+\mbox{}+ einfügt.
Ein entsprechendes Beispiel stellt die Überschrift dieses Abschnitts dar.
Dasselbe gilt für die ungünstig umbrochenen Wörter
(so empfiehlt der KSP Verlag darauf zu achten, dass keine seitenübergreifenden Wort-Trennungen und mehr als 4 nacheinanderfolgenden Zeilen mit Worttrennung am Ende der Zeile entstehen).
Weiteren Spielraum gibt das umbruchgeschützte Leerzeichen (eingefügt mit einer Tilde),
der Parameter \printkeyword{emergencystretch}, welcher global die maximal mögliche Dehnung des Zwischenraums zwischen den einzelnen Wörtern steuert,
sowie zusätzliche Steuermöglichkeiten der Silbertrennung, die auf Wikibooks%
\footnote{\url{https://de.wikibooks.org/wiki/LaTeX-W%C3%B6rterbuch:_Silbentrennung}}
beschrieben sind.